\documentclass[11pt]{article.cls}
\usepackage[utf8]{inputenc}
\usepackage{graphicx}
\usepackage{hyperref}
\usepackage{listings}
\usepackage{color}
\usepackage{float}

\title{Artificial Intelligence: Assignment 1 – Sokoban}
\author{Student Name}
\date{April 20, 2025}

\begin{document}
\maketitle

\section{Implementation and Optimizations}

\subsection{LRTA* (Learning Real-Time A*)}
We implemented LRTA* with the following Sokoban-specific optimizations:
\begin{itemize}
    \item State visit limitation to prevent cycling
    \item Heuristic value caching
    \item Differentiated costs for push vs. pull moves
    \item Backtracking mechanisms for deadlock situations
    \item Periodic clearing of visited states for memory management
\end{itemize}

\subsection{Simulated Annealing}
For Simulated Annealing, we implemented the following optimizations:
\begin{itemize}
    \item Adaptive cooling schedule based on solution progress
    \item Informed neighbor generation using heuristics
    \item Best-state memory retention
    \item Periodic reheating to escape local minima
\end{itemize}

\section{Implemented Heuristics}

\subsection{Basic Heuristic}
\begin{itemize}
    \item Manhattan distance between boxes and their nearest targets
    \item Penalty for boxes blocked against walls
\end{itemize}

\subsection{Advanced Heuristic}
\begin{itemize}
    \item Pattern database for common box configurations
    \item Deadlock analysis for invalid state avoidance
    \item Target accessibility evaluation
\end{itemize}

\section{Comparative Analysis}

\subsection{Execution Time}
LRTA* is faster on simple maps (easy_map1, easy_map2), while Simulated Annealing performs better on complex maps due to its stochastic nature.

\subsection{Number of Explored States}
LRTA* systematically explores the state space, while Simulated Annealing makes guided random jumps, resulting in less exhaustive but potentially more efficient exploration in large search spaces.

\subsection{Solution Quality}
\begin{itemize}
    \item LRTA*: Finds optimal or near-optimal solutions but may fail on complex maps
    \item Simulated Annealing: Sub-optimal but consistent solutions, with better adaptability to complex maps
\end{itemize}

\section{Conclusions}
Both algorithms present distinct advantages and disadvantages:
\begin{itemize}
    \item LRTA* is preferable for simple maps where finding the optimal solution is a priority
    \item Simulated Annealing is more suitable for complex maps where finding an acceptable solution is more important than optimality
\end{itemize}

\end{document}
